\chapter{Literature Review}
\label{ch:lit}

\section{Scope and method of the review}
\label{sec:lit-scope}

Five bodies of work bear on this project, and few papers cite across them. The psychology of emotion, emotion detection in text, image generation from text, the scoring of generated pictures, and the study of how people judge feeling in a picture each have something to say about a drawing made from the words ``tired but happy''. All five were read across to check that none of them had already asked this project's question. The starting point was the eight papers Dr Chen provided when the project began. From each of those the reference lists were followed outward, one literature at a time. Ninety-seven papers were read in full.

Every paper got two scores, one for how close it sat to the path of the dissertation and one for the strength of its evidence, and a paper was kept only when it did well on both. In practice that meant an empirical study with an evaluation of its own, or a dataset, lexicon or benchmark that later work builds on. Surveys and position pieces went out. So did careful work that was aimed somewhere else, and a reason was logged for every exclusion, so that the 24 papers dropped at this stage can be checked against the criterion. Two foundations would have failed the evidence test applied strictly, Ekman on basic emotions and Russell on the circumplex, and they were kept regardless, because every classifier and every lexicon further down assumes one or the other. That left 73 papers, and they fell into eight themes. On 9 September, after the second audit, 18 references were added for the models and the statistical procedures that the methodology depends on, which takes the bibliography to 91 entries; they are cited where they are used, here and in Chapter~\ref{ch:method}. The eight themes give this chapter its sections.

One term needs settling before the review starts. When other people's papers discuss how well an image corresponds to its prompt they say \emph{alignment}, and their term is kept here whenever their work is described. For this project the word is \emph{agreement}. What is measured here is how often two readers, one given the prompt and one given the drawing, return the same thing, and calling that alignment would imply a standard that the drawing either meets or misses, and this corpus has no such standard.

\section{Foundations: emotion theory and taxonomy}
\label{sec:lit-theory}

The vocabulary comes from two traditions. Ekman \cite{ekman1992argument} made the case that a small number of basic emotions exist, each with its own signal, and the six-way scheme that most classifiers still use is his: happiness, sadness, anger, fear, surprise and disgust. Russell's circumplex \cite{russell1980circumplex} puts affect on two continuous axes instead, valence and arousal, and the later paper by Posner, Russell and Peterson \cite{posner2005circumplex} argues from neural and developmental evidence that the dimensional picture fits better than the categorical one. Both are used here. The seven families in this dissertation, Ekman's six with neutral added, form the categorical rung, and the dimensional rung comes from where each family sits on the circumplex.

Neither scheme on its own covers what people report. Cowen and Keltner \cite{cowen2017self} asked people what they felt while watching more than two thousand brief video clips, and the reports resolved into 27 distinct categories linked by continuous gradients; the categories explained more of the variance than valence and arousal did, and neither description was enough by itself. The similarity framing rests on this result. If ``happy'' and ``excited'' lie next to each other on a gradient, then a rule that counts a drawing read as excited as a complete miss for a prompt that said happy is really measuring where the boundaries of the label set fall.

The same point has been made for faces in Dr Chen's own field. Using reverse correlation, Jack and colleagues showed that Western and East Asian observers do not hold the same mental representations of the six basic expressions \cite{jack2012facial}. A later study found that the signals in a dynamic expression arrive in a hierarchy over time, so that the early movements reduce six emotions to four, fear being confusable with surprise and disgust with anger at the start \cite{jack2014dynamic}. Dr Chen and her colleagues found distinct facial signatures for pain and for pleasure inside each culture \cite{chen2018distinct}. They then moved psychologically valid expressions onto social robots and a conversational virtual agent by measuring what observers reported seeing in a face \cite{chen2018reverse, chen2020building}. Their most recent work finds that facial movements code the category of an emotion and its intensity as separate components \cite{chen2024cultural}.

For a project that reads emotion off drawn faces, this has consequences. There is no reference drawing of an emotion to compare against, and near misses between neighbouring categories are to be expected. What a face signals is settled, in this tradition, by what human observers report, which this project does not measure because both of its readers are programs. So the foundations settled two design decisions, a small family scheme that two quite different readers can both output and a dimensional measure beside it, so that a drawing read as excited for the prompt ``happy'' earns partial credit.

\section{Emotion from text}
\label{sec:lit-text}

Text emotion recognition has mostly been built and tested on longer social-media text, and it is weakest on short first-person entries like the ones in this corpus. The reference resource is GoEmotions \cite{demszky2020goemotions}: 58,000 Reddit comments, each annotated with one or more of 27 fine-grained categories or with neutral, together with a mapping from those categories onto Ekman's families that the authors themselves propose and that this project uses. Their BERT (Bidirectional Encoder Representations from Transformers) model's macro-F1 was 0.46. The average hides a lot, because the score falls towards zero on the rare categories whose emotion is implied and never named. Harutyunyan and Kumar \cite{harutyunyan2026fine} re-ran the comparison on the same data across classical learners, a recurrent network and transformers, and the same shape appears in every row. Yang and colleagues \cite{yang2023context} used large language models to audit emotion datasets and found short texts under-specified as a rule; the label frequently turns on context that the text leaves out. A festival entry is that case, since it is written in the first person and is a few words long, and it often names two feelings.

Two papers take short texts on directly: Liu and colleagues \cite{liu2023emotion} propose an improved multi-label method for them, and Bhat \cite{bhat2024emotion} reports that a small transformer classifies short English texts well. Both were built and tested on tweets, though, and a tweet still gives the model more to work with than ``tired but happy'' does.

Dimensional resources go part of the way round the problem. EmoBank \cite{buechel2017emobank} rated ten thousand sentences on valence, arousal and dominance twice over, once from the writer's point of view and once from the reader's, and found that the two sets of ratings differ, which matters for this project because a visitor's words come from the writer's side and every reader in the pipeline sits, by construction, on the reader's side. The NRC lexicons supply valence, arousal and dominance norms at word level, first for twenty thousand English words \cite{mohammad2018obtaining} and now for fifty-five thousand \cite{mohammad2025nrc}, and the earlier crowdsourced lexicon \cite{mohammad2013crowdsourcing} maps words to basic emotions. Linguistic Inquiry and Word Count (LIWC) \cite{tausczik2010psychological} is still the transparent tool that psychologists reach for, a validated dictionary whose categories include positive and negative emotion, anxiety, anger and sadness. Zanwar and colleagues \cite{zanwar2022improving} found that fusing such psycholinguistic features with a transformer helps it generalise out of domain, which is a live worry for a corpus unlike anything the models were trained on. Niu and Mower Provost \cite{niu2025contrastive} compress what a large language model knows about emotion into a compact recogniser that needs no labelled examples. Parkington and colleagues \cite{parkington2025human} compared a large language model with expert human coders doing thematic analysis, and found high agreement between them but a loss of fine detail on emotional triggers. BERTopic \cite{grootendorst2022bertopic} was considered for pulling themes out of the corpus and set aside, because topic modelling has nothing to work with when the median document is four words long.

The rest of the text-side references are the models used in the pipeline. Entailment-based zero-shot classification \cite{yin2019benchmarking} turns each candidate label into a hypothesis sentence and puts it to a natural-language-inference model, which says whether the input entails it, and BART \cite{lewis2020bart} trained on MultiNLI \cite{williams2018broad} is the usual choice for that model and the one that reads the descriptions here. T5 \cite{raffel2020exploring} recast every task as text-to-text, its instruction-tuned successor \cite{chung2022scaling} extended that, and some of the vision--language models in Section~\ref{sec:lit-t2i} are built on the two of them. Everything runs through the Transformers library \cite{wolf2020transformers}.

On the text side the prompt reader is the least reliable part of the chain, and that shaped two choices. Its label on a four-word entry is uncertain, and every number downstream inherits the uncertainty. Chapter~\ref{ch:method} explains why GoEmotions was chosen over an entailment model and a hand-built lexicon for the prompt side; the short version is that it was trained on first-person Reddit comments, the nearest register to a festival entry available. The decision to keep the prompts that the reader labels neutral as a group of their own and analyse them, the specificity check of Chapter~\ref{ch:results}, comes from this literature too. A label of neutral on a two-word entry is exactly the kind of judgement Yang and colleagues warn about.

\section{Datasets, benchmarks and lexicons}
\label{sec:lit-data}

Every image-side resource in this literature is photographic or painted, and every one is in colour. EmoSet \cite{yang2023emoset} is the largest, at 118,000 images, each labelled with one of eight emotions and a set of attributes. Machajdik and Hanbury \cite{machajdik2010affective} built their affective classifier from colour, texture and composition features drawn from art theory, and they reported how far people agree with one another on the same images as a reference point. The Nencki Affective Picture System \cite{marchewka2014nencki} is a set of photographs that come with normative ratings of valence and arousal. EMOTIC \cite{kosti2020context} annotates people in their surroundings, with 26 categories as well as continuous dimensions, and its finding is that what surrounds a person changes which emotion annotators assign to that person. None of these resources contains line drawings, and that domain shift is one of the two reasons for the shape of the image side (the other, Dr Chen's direction, comes up in Section~\ref{sec:lit-t2i}). A classifier built on any of these resources could be run over the drawings, but its output could not be trusted, so the pipeline describes a drawing first and extracts the emotion from the description afterwards. The direct-CLIP reader is kept as the baseline that the method is compared against.

The art datasets come closest to this corpus because each pairs a picture with a short piece of affective text. ArtEmis \cite{achlioptas2021artemis} gathered 439,000 utterances in which annotators say what emotion an artwork makes them feel and why. ArtELingo \cite{mohamed2022artelingo} took the same paintings into four languages and found that the emotion labels move with the language and culture of the annotator. WikiArt Emotions \cite{mohammad2018wikiart} covers four thousand paintings and gave two results relied on later. Giving an image a title shifts the emotion people report for it. Images with a face in them get much more consistent emotion judgements from annotators than images without one. On the text side, CARER \cite{saravia2018carer} built a tweet corpus by distant supervision with eight emotions, and it is where the widely used ``emotion'' dataset comes from; the SemEval affect tasks from 2018 \cite{mohammad2018semeval} and from 2025 \cite{muhammad2025semeval} fix the shared benchmarks and the protocol for scoring intensity. Then there are the human-preference datasets built for text-to-image models, ImageReward \cite{xu2023imagereward}, Pick-a-Pic \cite{kirstain2023pick} and Human Preference Score \cite{wu2023human}, each of which records which of a pair of generated images people preferred, without asking about emotion, and trains a scorer on the answer.

\section{Text-to-image generative models}
\label{sec:lit-t2i}

The drawings were made by a DALL$\cdot$E model. The original DALL$\cdot$E \cite{ramesh2021zero} was a discrete autoregressive model, the second version \cite{ramesh2022hierarchical} decoded images by diffusion, conditioned on CLIP image embeddings, and DALL$\cdot$E~3 \cite{betker2023improving} trained on synthetic captions to follow prompts more closely and rewrites what the user typed before generating anything. Which of these made the festival drawings is not recorded. The provenance manifest inside each PNG names the software agent as ``DALL$\cdot$E'' and nothing more, with no version, no seed and no copy of the rewritten prompt. So no drawing can be produced again, and nothing about the generator by itself can be separated from the rest of the chain, a limit that returns in the consistency experiment. The diffusion models that DALL$\cdot$E sits beside, latent diffusion \cite{rombach2022high}, SDXL \cite{podell2023sdxl} and Imagen \cite{saharia2022photorealistic}, share a set of well-documented weaknesses: fine detail, rendered text and faces.

The readers used here are from the same few years. CLIP \cite{radford2021learning} learned to embed images and captions in one space from 400 million pairs of photographs and text, and it has become the default zero-shot image classifier. It was trained on colour photographs and its strength is objects, and for that reason Dr Chen's direction in late July was not to ask it for an emotion. Bootstrapping Language-Image Pre-training (BLIP) \cite{li2022blip} learned both captioning and visual question answering by bootstrapping from noisy web data; its VQA head takes an image and a free-form question and returns an answer, and the describer in this project does that five times per drawing. BLIP-2 \cite{li2023blip} joins a frozen image encoder to a large language model through a small trainable bridge. OpenCLIP \cite{ilharco2021openclip} is the open re-implementation of CLIP, and it is the model behind the direct-CLIP baseline. All of these readers are borrowed instruments, picked for what each is good at, and the choice to put a question-answering describer in front of the emotion reader, instead of a classifier, follows from the domain-shift argument of Section~\ref{sec:lit-data}.

\section{Image evaluation: metrics and frameworks}
\label{sec:lit-eval}

Automatic scoring of text-to-image models has settled on one question, whether the objects, attributes and relations that the prompt named are present in the picture. CLIPScore \cite{hessel2021clipscore} is the cosine similarity of the CLIP embeddings for the image and for the text. TIFA \cite{hu2023tifa} turns the prompt into questions and lets a VQA model answer them against the image. Cho and colleagues' Davidsonian Scene Graph \cite{cho2024davidsonian} orders those questions so that a dependent question is only asked once its premise has been confirmed. GenEval \cite{ghosh2023geneval} runs an object detector and checks counts, colours and positions. VQAScore \cite{lin2024evaluating, li2024evaluating} puts one yes/no question to a generative vision--language model, takes the probability of ``yes'' as the score, and on the human ratings of GenAI-Bench it tracks people far more closely than CLIPScore does. The Conditional Fr\'echet Distance \cite{koo2026evaluating} compares a model's output distribution with a reference distribution, conditioned on the prompt. Behind these sit the caption metrics that came first, SPICE \cite{anderson2016spice} and CIDEr \cite{vedantam2015cider}, and the Fr\'echet Inception Distance (FID) \cite{heusel2017gans}, which scores image quality without looking at the prompt.

How much any of this can be trusted is answered by two studies. Lee and colleagues \cite{lee2023holistic} evaluated text-to-image models on many axes at once and found that the automatic scores agree only weakly with human ratings on several of them, including photorealism and aesthetics. Zhang and Wan \cite{zhang2025rethinking} put CLIPScore, VQAScore and DSGScore under stress on a set of about a thousand prompts, which is roughly the size of this corpus, and the first and the third reversed the ranking of models when the seed changed and moved materially after a single pixel was perturbed. VQAScore was steadier, though not immune.

A second group of papers looks at the judge instead of the score. Chang and colleagues \cite{chang2023prompt} describe the way experienced prompt artists push a model towards a style with keywords such as ``black and white'' and ``line art'', and note that an output drifts towards the model's average look unless it is pushed. Oppenlaender and colleagues \cite{oppenlaender2024artworks} report that driving generation from a body pose changes the feeling of the resulting image without anyone intending it, and that uncanny artefacts change how the image is read emotionally. Carrasco and colleagues \cite{carrasco2024level} asked how closely human raters' labels follow the emotion a generator had been told to produce, scored it with Cohen's kappa and Krippendorff's alpha, and found the agreement much stronger once the labels were reduced to positive and negative. Across ten experiments, Horton, White and Iyengar \cite{horton2023bias} together with Millet and colleagues \cite{millet2023defending} find that people rate the same piece of work lower when they are told a machine made it, which constrains the design of any later human study on these drawings. Two benchmarks test the machine judges themselves. Bhattacharyya and Wang \cite{bhattacharyya2025evaluating} found that vision--language models do a poor job of recognising, zero-shot, the emotion an image evokes in a viewer, lean systematically towards positive labels, and change their answer with the order in which the candidate labels are listed. Gao and colleagues \cite{gao2025eemo} report that multimodal models manage valence, fail on arousal and dominance, and run together what the person in the picture expresses and what the person looking at it feels.

No alignment metric measures affect, so none of them could be the instrument here, although VQAScore, which works by asking a question and reading the answer, is the nearest relative of the describer used here. They are also unstable. Because the scores move with the seed and with a single pixel, no value in this dissertation is reported bare, and each is compared with its random-pairing null. The positive lean and the expressed-versus-evoked confusion that the two benchmarks found in machine judges are the failure modes the audits in Chapter~\ref{ch:results} check for, and both appear in this corpus.

\section{The intersection: emotion in generated images}
\label{sec:lit-intersection}

The papers that put emotion into generated images condition on a target and work in colour. EmoGen \cite{yang2024emogen} generates an image from an emotion category by way of an emotion space mapped onto CLIP's, and scores the result with an EmoSet-trained classifier. Make Me Happier \cite{lin2024make} takes an existing image and edits it towards a chosen emotion; its authors found positive emotions harder to separate from each other than negative ones, and fine facial detail was a weak point. EmotiCrafter \cite{dang2025emoticrafter} conditions generation on a continuous valence--arousal coordinate. Mehta and Buntain \cite{mehta2024generating} asked whether generators drift towards fear whatever the prompt says, with an image-emotion classifier for the pictures and a text-emotion classifier for the prompts, and that cross-modal design is the nearest methodological relative of the pipeline here. Zhou and colleagues \cite{zhou2025analyzing} turn an image into affective captions first and classify the evoked emotion from those, so the describe-then-extract idea has a precedent on the recognition side too.

Alvarado \cite{alvarado2025texttoimage} had a diffusion model render emotion words and then extracted colour features from what came back, and found that lightness and chroma track emotional meaning in the output more closely than hue does. Only colour features are modelled. Shape and line are not tested, so the study says nothing either way about what a drawing with no colour in it can do. The literature therefore offered no prediction for a corpus without colour, which is one reason it seemed worth studying.

The nearest prior study is by Lomas and colleagues \cite{lomas2024improved}. Twenty-four raters scored 240 images, made by three generators from a fixed list of ten emotions, with each emotion rendered on a person or on a robot, rating alignment from 0 to 10. The generating model was the largest effect, with DALL$\cdot$E~3 well ahead of Stable Diffusion and of DALL$\cdot$E~2, and the scores differed sharply from one emotion to the next; emotions rendered on a person scored higher than the same emotions on a robot. Their prompts were built from a template and their images were in colour. Nobody, in the papers reviewed, has measured affect agreement for self-reports that came straight from the public, mixed feelings included, and were drawn without colour. None of these studies drew the same words twice either, so none could ask how stable a reading is, and the repeated prompts in this corpus make that question possible.

\section{Emotion perception and human judgement}
\label{sec:lit-perception}

Machine recognition of emotion in pictures was built on photographs, and it mixes up two different constructs. Yang and colleagues \cite{yang2018weakly} learned visual sentiment with a pair of coupled networks and weak supervision. EmoSet and the other datasets of Section~\ref{sec:lit-data} supply the labels such systems learn from, and the benchmarks in Section~\ref{sec:lit-eval} show how brittle the resulting judges are. Underneath this is a distinction, formalised by Gao and colleagues \cite{gao2025eemo} and made concrete in Dr Chen's experiments, between what the person in a picture expresses and what a viewer feels on looking at it. A classifier fitted to crowd labels returns a mixture of the two. The pipeline here reads expressed cues only. It works from a drawn expression, a posture and an action. A study with human raters would be reading evoked affect, and the two cannot be pooled. The reverse-correlation methods that Jack and Dr Chen have used \cite{jack2012facial, chen2018distinct, chen2020building, chen2024cultural} show how such a study should be run: start by measuring what observers report seeing, without presupposing it.

The distinction settles the vocabulary used from here on. What the image reader returns is called a reading throughout. Nothing in the computational results says what a person sees, because nobody looked at a drawing until the single-coder check of Chapter~\ref{ch:results}, and that check is the only point in the project at which somebody does.

\section{Context}
\label{sec:lit-context}

The drawings were made at a public festival. One study of how the public receives machine-made art bears on them directly. In a study of 565 participants, Ragot and colleagues \cite{ragot2020ai} found that an identical artwork is rated lower as soon as it is labelled as made by a machine. Every drawing here was shown to its author as machine-made, so anyone who rates them in future will have to hide that label or counterbalance it.

\section{The gap}
\label{sec:lit-gap}

Taken together, these literatures describe an object that none of them has studied. A short self-report holds blended, graded affect, and short first-person text is where emotion classifiers do worst. A generator renders such a prompt fluently but can falter on faces and fine detail, and it may not even draw the words the visitor typed. The strongest evaluation metrics count objects and are unstable in their own right, while the work on emotion in images reads affect mainly through colour, on photographs and paintings, using judges that lean positive and confuse the emotion a depicted person expresses with the emotion felt by the viewer.

Across the 73 papers, no study measures how far the emotion a reader recovers from a first-person self-report that is short, unmediated and mixed in affect agrees with the emotion a reader recovers from a colourless line drawing generated from it. No study asks, either, how stable that second reading is when the identical words are drawn a second time. A compositional metric can pass a line drawing on the question ``is there a person?'' and the reader can still return sadness for a prompt that read ``happy''. Lomas and colleagues used a fixed typology and coloured images, whereas the prompts here are whatever somebody typed, blends included, and the drawings are colourless. The conditional generators treat emotion as a control setting; here it is latent in a visitor's words. The visual-emotion recognisers were trained on photographs and paintings, so monochrome line art sits outside their training distribution, and not one of them keeps expressed and evoked affect apart for such stimuli.

The rest of the dissertation asks two questions. The first is how closely two automatic readers agree about the emotion in a visitor's words and in the drawing made from them, measured against chance, and how much of that agreement depends on whether the drawing has a readable face. The second is how stable the reading stays when one prompt is drawn several times. What a person looking at these drawings would see needs people to answer it and is left to further work. The dissertation's contribution is the corpus and the question, a similarity ladder with a chance level under every rung, a consistency experiment that the repeated prompts make possible, and an audited toolkit; it proposes no new metric or classifier.
